\chapter{Protocollo seriale del LD2410B}
\label{app:protocollo}

Riferimento: \emph{LD2410B Serial communication protocol}
V1.07~\cite{hilinkProtoV107} --- sigla del frontespizio; il contenuto distribuito
corrisponde alla revisione 1.08 del 22 novembre 2024, per le ragioni discusse in
\cref{sec:protocollo-seriale}. Questa appendice riporta gli elementi effettivamente
usati dal firmware realizzato, per rendere il lavoro riproducibile senza dover
ricostruire il protocollo dal documento originale.

\section{Parametri della porta seriale}

\begin{center}\small
\begin{tabular}{@{}ll@{}}
\toprule
Baud rate & 256000 (fisso, non modificabile senza comando \texttt{0x00A1}) \\
Formato & 8 bit di dati, nessuna parità, 1 bit di stop (8N1) \\
Livelli & 3,3~V TTL \\
\bottomrule
\end{tabular}
\end{center}

Il baud di 256000 richiede una UART hardware: sull'ESP32 si usa \texttt{Serial2}
(GPIO25 come RX, GPIO26 come TX nella nostra configurazione). Un'implementazione a
software non regge questa velocità.

\section{Struttura dei frame}

\subsection{Frame di comando (dall'host al modulo)}

\begin{center}\small
\begin{tabular}{@{}llp{0.42\linewidth}@{}}
\toprule
\textbf{Campo} & \textbf{Lunghezza} & \textbf{Valore} \\
\midrule
Intestazione & 4 byte & \texttt{FD FC FB FA} \\
Lunghezza dati & 2 byte & little-endian \\
Parola di comando & 2 byte & little-endian \\
Valore del comando & variabile & dipende dal comando \\
Chiusura & 4 byte & \texttt{04 03 02 01} \\
\bottomrule
\end{tabular}
\end{center}

Il modulo risponde con un frame della stessa struttura, in cui la parola di comando
ha il bit più significativo posto a uno, seguita da uno stato (zero = successo) e
dagli eventuali dati di risposta.

\subsection{Frame di dati (dal modulo all'host)}

\begin{center}\small
\begin{tabular}{@{}llp{0.42\linewidth}@{}}
\toprule
\textbf{Campo} & \textbf{Lunghezza} & \textbf{Valore} \\
\midrule
Intestazione & 4 byte & \texttt{F4 F3 F2 F1} \\
Lunghezza dati & 2 byte & little-endian \\
Tipo di dato & 1 byte & \texttt{0x01} = engineering mode, \texttt{0x02} = normale \\
Testa & 1 byte & \texttt{0xAA} \\
Stato del bersaglio & 1 byte & maschera: \texttt{0x01} in movimento, \texttt{0x02} stazionario \\
Distanza in movimento & 2 byte & cm, little-endian \\
Energia in movimento & 1 byte & 0--100 \\
Distanza stazionaria & 2 byte & cm, little-endian \\
Energia stazionaria & 1 byte & 0--100 \\
Distanza di rilevamento & 2 byte & cm, little-endian \\
\multicolumn{3}{@{}l@{}}{\emph{solo in engineering mode:}} \\
Gate massimo movimento & 1 byte & \\
Gate massimo stazionario & 1 byte & \\
Energie per-gate movimento & 9 byte & gate 0--8, valori 0--100 \\
Energie per-gate stazionario & 9 byte & gate 0--8, valori 0--100 \\
Sensore di luce & 1 byte & 0--255 \\
Stato del pin OUT & 1 byte & \\
\multicolumn{3}{@{}l@{}}{} \\
Coda & 1 byte & \texttt{0x55} \\
Controllo & 1 byte & \texttt{0x00} \\
Chiusura & 4 byte & \texttt{F8 F7 F6 F5} \\
\bottomrule
\end{tabular}
\end{center}

\section{Comandi utilizzati}

\begin{center}\small
\begin{tabularx}{\linewidth}{@{}llX@{}}
\toprule
\textbf{Parola} & \textbf{Funzione} & \textbf{Note d'uso} \\
\midrule
\texttt{0x00FF} & entra in configurazione & obbligatorio prima di ogni altro
  comando; valore \texttt{0x0001} \\
\texttt{0x00FE} & esce dalla configurazione & il modulo riprende a inviare frame di
  dati \\
\texttt{0x0062} & attiva engineering mode & va inviato dentro la modalità
  configurazione \\
\texttt{0x0063} & disattiva engineering mode & \\
\texttt{0x00A0} & legge la versione firmware & restituisce tipo e versione
  (nel nostro caso 2.44.25070917) \\
\texttt{0x00A5} & legge il MAC Bluetooth & \\
\texttt{0x0061} & legge i parametri correnti & restituisce risoluzione, gate massimi,
  timeout e le 9+9 soglie; \textbf{è la lettura da usare per i parametri}, non lo
  strumento da PC \\
\texttt{0x0060} & imposta i parametri & usato per limitare il gate massimo nella prova
  di selettività spaziale \\
\texttt{0x00AA} & imposta la risoluzione di distanza & 0,75~m oppure 0,20~m per gate \\
\texttt{0x000B} & auto-calibrazione del rumore di fondo & disponibile dal firmware
  V2.44~\cite{hilinkV244} \\
\bottomrule
\end{tabularx}
\end{center}

\section{Regola di rilevamento}

Dal protocollo V1.07, \S1.2.2: il bersaglio in un gate è riconosciuto
\textbf{solo se l'energia del gate supera la soglia configurata per quel gate}. Da
questa regola derivano due fatti osservati sperimentalmente:

\begin{itemize}[leftmargin=1.4em,itemsep=2pt]
  \item i canali stazionari dei gate 0 e 1, che hanno soglia di fabbrica pari a~0 e
        energia riportata sempre nulla, sono disattivati per costruzione
        (\cref{sec:senergy-zero});
  \item le soglie influiscono sulla \emph{decisione} di rilevamento, non sul valore
        di energia riportato: non possono quindi risolvere la saturazione
        (\cref{sec:saturazione}).
\end{itemize}

\section{Auto-calibrazione del firmware V2.44}

Il documento delle novità del firmware V2.44~\cite{hilinkV244} descrive il
rilevamento automatico del rumore di fondo, che tara le soglie di movimento e
stazionarie. La procedura richiede 10~s per uscire dal campo del sensore più 60~s di
misura, per un totale di 70~s da trascorrere fuori dalla portata del modulo, e
l'applicazione mobile in versione adeguata (Android $\geq$ V1.5.12 oppure
iOS $\geq$ V1.5.4).

Due precisazioni utili, entrambe verificate:

\begin{itemize}[leftmargin=1.4em,itemsep=2pt]
  \item la funzione ``Detect noise floor'' presente nella schermata dei parametri
        dell'applicazione è \emph{soltanto} una funzione di visualizzazione: mostra i
        valori di rumore misurati ma non li applica come soglie. Non va confusa con
        l'auto-calibrazione;
  \item il documento non descrive alcuna procedura di aggiornamento del firmware, e
        né lo strumento da PC né l'applicazione mobile offrono una funzione di flash.
        La versione del firmware è quindi da considerare un dato dell'esemplare.
\end{itemize}

\section{Il frame binario del \ldc\ (modalità non documentata)}
\label{app:ld2420-frame}

Il produttore documenta per il \ldc\ soltanto i comandi di configurazione
(\cite{hilinkProtoLD2420}: apertura e chiusura della modalità comandi con
\texttt{0xFF} e \texttt{0xFE}, versione \texttt{0x00}, riavvio \texttt{0x68},
lettura e scrittura dei parametri \texttt{0x08} e \texttt{0x07}) e nessun frame di
uscita dei dati. La modalità usata in questa tesi è ricostruita dal componente
\texttt{ld2420} di ESPHome~\cite{esphomeLD2420src} e va considerata una scelta
implementativa di comunità, non una specifica.

\begin{itemize}[leftmargin=1.4em,itemsep=2pt]
  \item \textbf{commutazione}: in modalità comandi, \texttt{0x0012} con parametro
        \texttt{0x0000} e valore \texttt{0x0004} (modalità \emph{energy}) oppure
        \texttt{0x0064} (modalità \emph{simple}, l'ASCII di fabbrica);
  \item \textbf{frame di dati}, 45 byte: intestazione \texttt{F4 F3 F2 F1},
        lunghezza (2 byte), presenza (1 byte, offset 6), distanza in centimetri
        (\texttt{uint16} little-endian, offset 7), sedici energie per gate
        (\texttt{uint16} little-endian, offset 9--40), chiusura
        \texttt{F8 F7 F6 F5}. Le intestazioni sono le stesse del \ldb;
  \item \textbf{cadenza} 10~Hz; energie intere 0--65535, mostrate dallo strumento del
        produttore come $10\log_{10}$ del valore grezzo, stessa relazione delle
        soglie;
  \item \textbf{parametri} (comando \texttt{0x07}): \texttt{0x0000} gate minimo,
        \texttt{0x0001} gate massimo, \texttt{0x0004} ritardo di scomparsa in
        secondi, \texttt{0x10}--\texttt{0x1F} soglie di attivazione per gate,
        \texttt{0x20}--\texttt{0x2F} soglie di mantenimento. La scrittura via UART
        vive in RAM e si perde al ciclo di alimentazione del modulo.
\end{itemize}

Sull'esemplare in dotazione la scrittura di un gate minimo non modifica la decisione
di presenza e confina la distanza riportata alla finestra ammessa; il gate massimo
scritto via UART è invece efficace sulla decisione (\cref{sec:ld2420-limiti}).
